\slajdid{09_2-s057}
\begin{frame}[label=09_2-s057]
\gusttekst\setlength{\abovedisplayskip}{3pt}\setlength{\belowdisplayskip}{3pt}\setlength{\abovedisplayshortskip}{2pt}\setlength{\belowdisplayshortskip}{2pt}\[t_p=0.69R_{eq}C_{eq}\]
Da bi obezbedili da je ovo kašnjenje manje ili treba smanjiti $R_{eq}$ ili $C_{eq}$.
\[R_p\approx\frac34\frac{V_{DD}}{I_{Dpsat}}\left(1-\frac79|\lambda_p|V_{DD}\right)\]
i
\[I_{Dpsat}=\frac{k_p}2\frac1{1+\frac{-V_{DD}-V_{Tp}}{E_{Cp}L_p}}(-V_{DD}-V_{Tp})^2\]
Sigurno važi $R_{eq}\sim\frac1{V_{DD}}$, odnosno da bi smanjili dinamičke otpornosti a time i kašnjenje moramo povećati napon napajanja. Na žalost tada ćemo disipaciju povećati sa trećim stepenom pa moramo obezbediti znatno bolje hlađenje procesora. Overklokovanje bez povećanja napona napajanja neće uspeti, a bez povećanja hlađenja dovešće do stradanja procesora. Ali isto tako nam ovi rezultati pokazuju da je u sistemima koji su baterijski napajani, ograničen izvor energije, moguće uštedeti energiju, produžiti rad uređaja, tako što će se procesoru i digitalnim sistemima kada nisu potrebni smanjivati učestanost rada i napon napajanja. Tehnika (voltage frequency scaling) koja se koristi u svim mobilnim telefonima, laptopovima itd…
\end{frame}
